% gerado por figuras_artigo.py -- não editar à mão
\begin{table}[t]
  \centering
  \caption{Segmentação no split \emph{test}, camada \texttt{chord[0]}. \emph{overseg}, \emph{underseg} e \emph{seg} em \% (média ponderada pela duração; \emph{seg} é o mínimo dos dois primeiros). \emph{Segm. est.} e \emph{Segm. ref.} são o número médio de segmentos por faixa, estimado e de referência (média simples). Fonte: \texttt{ablation\_20260927\_014253\_ref0.csv}.}
  \label{tab:segmentacao}
  \footnotesize
  \begin{tabular}{lrrrrr}
    \toprule
    Variante & overseg & underseg & seg & Segm. est. & Segm. ref. \\
    \midrule
    \texttt{V4} & 75,51 & 88,13 & 74,14 & 62,7 & 12,0 \\
    \texttt{C} & 59,41 & 88,50 & 58,69 & 30,5 & 12,0 \\
    \texttt{D} & 74,62 & 83,15 & 70,24 & 19,8 & 12,0 \\
    \texttt{E} & 84,15 & 71,89 & 68,77 & 10,8 & 12,0 \\
    \texttt{F0rg} & 73,57 & 88,20 & 72,66 & 20,9 & 12,0 \\
    \texttt{D0rg} & 71,29 & 87,35 & 68,81 & 22,5 & 12,0 \\
    \bottomrule
  \end{tabular}
\end{table}
